\documentclass[11pt]{article}
\usepackage[margin=1in]{geometry}
\usepackage[T1]{fontenc}
\usepackage[utf8]{inputenc}
\usepackage{hyperref}
\hypersetup{colorlinks=true,urlcolor=blue,linkcolor=black}
\pagestyle{empty}
\setlength{\parskip}{0.4em}
\setlength{\parindent}{0pt}
\begin{document}

\begin{flushleft}
Ke Cheng Hsu \\
Department of Economics \\
The Hong Kong University of Science and Technology \\
Clear Water Bay, Kowloon, Hong Kong \\
\href{mailto:khsu@connect.ust.hk}{khsu@connect.ust.hk} \\
\href{https://cliohsu1997.github.io/}{https://cliohsu1997.github.io/}
\end{flushleft}

\vspace{0.4em}
\today

\vspace{0.6em}
Faculty Search Committee \\
Social Sciences / Economics \\
Duke Kunshan University \\
Kunshan, Jiangsu, China

\vspace{0.6em}
Dear Members of the Search Committee,

I am writing to apply for the faculty position in Economics on the Professor of the Practice teaching track at Duke Kunshan University (Academic Jobs Online \#32544). I am a Ph.D.\ candidate in Economics at the Hong Kong University of Science and Technology (expected June 2027). My field is empirical industrial organization and applied microeconomics.

I am particularly drawn to DKU's research-inflected liberal arts curriculum, where interdisciplinary communities take the place of departmental majors and faculty teach in clusters across fields. This structure matches how I teach: start from a concrete market question, build intuition with an example, then develop the formal model. The seven-week modular terms and discussion-based classes of twenty or fewer suit this style, and the two-year Signature Work program fits my commitment to close, tailored supervision of student projects. At HKUST I already enjoy working with people from different cultures and backgrounds, and DKU's international community---students from more than seventy countries earning both Duke and DKU degrees---is where I want to build my career on the teaching track.

At HKUST I teach and assist in undergraduate and postgraduate economics courses in English, including microeconomic theory, information economics, and macroeconomics; for some courses I also designed tutorial materials and exam questions. My average teaching-assistant evaluation is 4.72/5. As lecturer for Math Camp for the MSc in Economics, I teach linear algebra, statistics, dynamic programming, real analysis, and optimization, with an instructor evaluation of 4.53/5. I have also served as Teaching Assistant Coordinator in the Business School, coaching postgraduate instructors on materials and practice teaching. I am prepared to teach principles and intermediate economics and, at the graduate level, Ph.D.\ microeconomic theory, information economics, applied microeconometrics, and empirical industrial organization---and to contribute to advising and program service.

My research on return policy, consumer search, and pricing and inventory under resale price maintenance gives me classroom examples that connect theory to markets students recognize. A separate statement describes how I integrate teaching and research. Thank you for your consideration.

\vspace{1em}
Sincerely,

\vspace{1em}
Ke Cheng Hsu \\
Ph.D.\ Candidate, Department of Economics, HKUST

\end{document}
